\subsection{Local-to-global mathematics: sheaves, closure systems and why the names matter}

Several mathematical traditions offer precise versions of ``local pieces forming a global object,'' and their differences help prevent metaphor from becoming false formalism. Sheaf theory studies compatible local data and the conditions under which local sections can be glued into global sections. In contextuality, the absence of a global section can itself carry substantive meaning \cite{abramsky2011}. Recent AI-science work uses sheaf-theoretic transport and obstruction to diagnose whether a representational language can coherently extend into a new regime \cite{olivieri2026sheaf}. This is close to \RAKL{} in vocabulary and motivation. It narrows the novelty claim and supplies a useful comparator: the cited framework is a finite diagnostic of transport and theory-shift candidates, while \RAKL{} embeds obstruction inside a larger evidence, authority, memory, search and method-evolution control plane.

Formal Concept Analysis (FCA) offers a different route to genuine lattice structure. Given a formal context relating objects to attributes, FCA constructs formal concepts and orders them into a concept lattice \cite{ganterwille2024}. This matters because the historical name ``Recursive Atomic Knowledge Lattices'' can be read as an order-theoretic claim. The current atom--witness--path implementation does not obtain a lattice merely because atoms are combinable. Its honest generic structure is a typed compatibility complex. A local order-theoretic lattice can be licensed only after a scoped order or closure operator is defined and the relevant meet/join obligations are established. FCA is therefore a useful mathematical neighbor and a terminology constraint, not a hidden implementation detail.

The distinction can be summarized as follows. A graph gives adjacency. A compatibility complex gives witnessed co-admissibility among typed objects and constructive paths. A poset gives an antisymmetric order. A lattice additionally supplies meets and joins for the relevant pairs. A closure system supplies a family of closed objects and can induce lattice structure under the appropriate conditions. A sheaf organizes local data, restrictions and gluing. These structures can coexist in one research architecture, but none is interchangeable with another. \RAKL{} uses different ones for different jobs: provenance is graph-like, authority is partially ordered, local chart synthesis is atlas/sheaf-like when the obligations are defined, and constructive candidate assembly currently uses a compatibility complex.

This layered view also changes what ``global'' means. A global model may fail to exist because local charts genuinely conflict; it may exist but be non-unique because the evidence underdetermines it; or it may be representable but scientifically unsupported. Global existence, uniqueness and authority are therefore three separate predicates. A successful gluing calculation cannot by itself increase evidence quality, and strong evidence for local charts cannot guarantee a unique global synthesis.
